\documentclass[12pt,a4paper]{article}
\usepackage[utf8]{vietnam}
\usepackage{amsmath,amssymb}
\usepackage{tabularx}
\usepackage{tikz}
\usetikzlibrary{calc,angles,quotes,arrows.meta,patterns}
\usepackage{xcolor}
\usepackage{enumitem}
\usepackage{fancyhdr}
\usepackage{graphicx}
\usepackage[top=0.9cm, bottom=0.9cm, left=1.2cm, right=1.2cm]{geometry}
\usepackage[dethi]{ex_test}
\SAOPTN{widthSA=3.2,heightSA=0.62}

\setlength{\headheight}{15pt}
\setlength{\parskip}{0pt}

\pagestyle{fancy}
\fancyhf{}
\rhead{\textit{Môn Toán 9 -- Đề số 05}}
\lhead{\textit{Bộ đề ôn tập Giữa học kì I}}
\cfoot{\thepage}

\begin{document}

% --- KHUNG TIÊU ĐỀ ĐỀ THI ---
\noindent
\begin{tabularx}{\linewidth}{X r}
\textbf{TRƯỜNG THCS: \dotfill} & \textbf{ĐỀ KIỂM TRA ĐỊNH KỲ GIỮA HỌC KỲ I} \\
\textbf{Họ và tên: \dotfill} & \textbf{NĂM HỌC 2026 -- 2027} \\
\textbf{Lớp: \dotfill \quad SBD: \dotfill} & \textbf{Môn: TOÁN 9} \\
\hline
\end{tabularx}
\vspace{2mm}

\begin{center}
  \textbf{\large ĐỀ ÔN TẬP GIỮA HỌC KỲ I -- ĐỀ SỐ 05}\\
  \textit{Thời gian làm bài: 90 phút (Không kể thời gian phát đề)}
\end{center}
\vspace{1mm}

% --- PHẦN I ---
\noindent\textbf{\large PHẦN I. CÂU TRẮC NGHIỆM NHIỀU PHƯƠNG ÁN LỰA CHỌN (3,0 điểm)}

\textit{Thí sinh trả lời từ câu 1 đến câu 12. Mỗi câu hỏi thí sinh chỉ chọn một phương án duy nhất.}

\vspace{1mm}
\begin{ex}
Cho hệ phương trình $\begin{cases} 3x - y = 2 \\ x + y = 6. \end{cases}$ Nghiệm $(x; y)$ của hệ phương trình là:
\choice
{\True $(2; 4)$}
{$(4; 2)$}
{$(1; 5)$}
{$(3; 3)$}
\end{ex}

\begin{ex}
Hệ phương trình $\begin{cases} 2x - y = 3 \\ 4x - 2y = 6 \end{cases}$ có số nghiệm là:
\choice
{Vô nghiệm}
{Nghiệm duy nhất}
{\True Vô số nghiệm}
{Hai nghiệm phân biệt}
\end{ex}

\begin{ex}
Cho $a < b$. Khẳng định nào sau đây là đúng?
\choice
{-2a < -2b}
{\True $a - 5 < b - 5$}
{$a + 3 > b + 3$}
{$-3a < -3b$}
\end{ex}

\begin{ex}
Nghiệm của bất phương trình $3x - 6 \ge 0$ là:
\choice
{$x \le 2$}
{$x < 2$}
{\True $x \ge 2$}
{$x > 2$}
\end{ex}

\begin{ex}
Cho tam giác $ABC$ vuông tại $A$. Biểu thức nào sau đây tính $\cos B$?
\choice
{$\frac{AC}{BC}$}
{\True $\frac{AB}{BC}$}
{$\frac{AC}{AB}$}
{$\frac{AB}{AC}$}
\end{ex}

\begin{ex}
Giá trị của biểu thức $M = \sin^2 25^\circ + \sin^2 65^\circ - \tan 40^\circ \cdot \cot 40^\circ$ là:
\choice
{\True $0$}
{$1$}
{$-1$}
{$2$}
\end{ex}

\begin{ex}
Mua $3$ quyển vở và $2$ cái bút bi hết $31\,000\text{ đồng}$. Mua $2$ quyển vở và $1$ cái bút bi cùng loại hết $20\,000\text{ đồng}$. Giá tiền của một quyển vở và một cái bút bi lần lượt là:
\choice
{$8\,000\text{ đồng}$ và $3\,000\text{ đồng}$}
{\True $9\,000\text{ đồng}$ và $2\,000\text{ đồng}$}
{$10\,000\text{ đồng}$ và $1\,000\text{ đồng}$}
{$7\,000\text{ đồng}$ và $5\,000\text{ đồng}$}
\end{ex}

\begin{ex}
Cho tam giác $ABC$ vuông tại $A$, đường cao $AH$. Biết $BH = 4\text{ cm}, CH = 9\text{ cm}$. Độ dài cạnh $AB$ bằng:
\choice
{\True $2\sqrt{13}\text{ cm}$}
{$6\text{ cm}$}
{$5\text{ cm}$}
{$3\sqrt{13}\text{ cm}$}
\end{ex}

\begin{ex}
Cho tam giác $ABC$ vuông tại $A$, cạnh $AC = 6\text{ cm}$, góc $\widehat{C} = 30^\circ$. Độ dài cạnh huyền $BC$ là:
\choice
{$3\text{ cm}$}
{$12\text{ cm}$}
{$6\sqrt{3}\text{ cm}$}
{\True $4\sqrt{3}\text{ cm}$}
\end{ex}

\begin{ex}
Điều kiện xác định của phương trình $\frac{x+1}{x-3} = \frac{18}{x^2-9} + 3$ là:
\choice
{$x \neq 3$}
{$x \neq -3$}
{\True $x \neq \pm 3$}
{$x \neq 3$ hoặc $x \neq -3$}
\end{ex}

\begin{ex}
Tìm giá trị của tham số $m$ để hệ phương trình $\begin{cases} mx + y = 3 \\ x + my = 3 \end{cases}$ có vô số nghiệm:
\choice
{\True $m = 1$}
{$m = -1$}
{$m = 2$}
{$m = 0$}
\end{ex}

\begin{ex}
Một cái thang dài $5\text{ m}$ dựa vào tường, góc tạo bởi thang và mặt đất là $65^\circ$. Khoảng cách từ chân thang đến chân tường (làm tròn đến hai chữ số thập phân) là:
\choice
{$4{,}53\text{ m}$}
{\True $2{,}11\text{ m}$}
{$2{,}54\text{ m}$}
{$5{,}52\text{ m}$}
\end{ex}


\vspace{1mm}
\noindent\textbf{BẢNG TRẢ LỜI TRẮC NGHIỆM PHẦN I:}
\begin{center}
\begin{tabular}{|c|c|c|c|c|c|c|c|c|c|c|c|c|}
\hline
\textbf{Câu} & \textbf{1} & \textbf{2} & \textbf{3} & \textbf{4} & \textbf{5} & \textbf{6} & \textbf{7} & \textbf{8} & \textbf{9} & \textbf{10} & \textbf{11} & \textbf{12} \\
\hline
\textbf{Chọn} & \rule{0pt}{0.65cm} & & & & & & & & & & & \\
\hline
\end{tabular}
\end{center}
\vspace{2mm}

\setcounter{ex}{0}
\noindent\textbf{\large PHẦN II. CÂU TRẮC NGHIỆM ĐÚNG SAI (2,0 điểm)}

\textit{Thí sinh trả lời câu 1 và câu 2. Trong mỗi ý a), b), c), d) ở mỗi câu, thí sinh chọn đúng hoặc sai.}

\vspace{1mm}
\begin{ex}
Cho bất phương trình $(3m - 1)x < 2$. Xét tính Đúng/Sai của các khẳng định sau:
\choiceTF
{\True Khi $m = 0$, bất phương trình trở thành $-x < 2$ và có nghiệm $x > -2$}
{\True Khi $m = 1$, bất phương trình có nghiệm là $x < 1$}
{Khi $m = -1$, bất phương trình có nghiệm là $x < -\frac{1}{2}$}
{\True Khi $m = -2$, bất phương trình có nghiệm nguyên nhỏ nhất là $0$}
\end{ex}

\begin{ex}
Cho hệ phương trình $\begin{cases} 4x - y = -5 \\ x + y = 5. \end{cases}$ Xét tính Đúng/Sai của các khẳng định sau:
\choiceTF
{\True Hệ phương trình trên là hệ hai phương trình bậc nhất hai ẩn}
{\True Cộng từng vế hai phương trình của hệ ta được phương trình $5x = 0$}
{\True Nghiệm của hệ phương trình là cặp số $(0; 5)$}
{Giá trị của biểu thức $P = 2x + 3y$ tại nghiệm của hệ bằng $12$}
\end{ex}


\noindent\textbf{BẢNG TRẢ LỜI PHẦN II:}
\begin{center}
\begin{tabular}{|c|c|c|c|c||c|c|c|c|c|}
\hline
\textbf{Câu 1} & \textbf{a)} & \textbf{b)} & \textbf{c)} & \textbf{d)} & \textbf{Câu 2} & \textbf{a)} & \textbf{b)} & \textbf{c)} & \textbf{d)} \\
\hline
\textbf{Đ/S} & \rule{0pt}{0.6cm} & & & & \textbf{Đ/S} & & & & \\
\hline
\end{tabular}
\end{center}
\vspace{2mm}

\setcounter{ex}{0}
\noindent\textbf{\large PHẦN III. CÂU HỎI TRẢ LỜI NGẮN (2,0 điểm)}

\textit{Thí sinh trả lời câu 1 và câu 2. Điền kết quả vào ô trống tương ứng.}

\vspace{1mm}
\begin{ex}
Một xe tải chở sữa tươi, mỗi thùng sữa nặng trung bình $10\text{ kg}$. Trọng tải tối đa của xe là $5{,}25\text{ tấn}$. Bác lái xe nặng $65\text{ kg}$. Hỏi xe có thể chở được tối đa bao nhiêu thùng sữa như vậy? \shortans{$518$}
\end{ex}

\begin{ex}
Với giá trị nào của $x$ thì hai biểu thức $A = \frac{3}{3x+1} + \frac{2}{1-3x}$ và $B = \frac{x-5}{9x^2-1}$ có cùng một giá trị? \shortans{$x = 0$}
\end{ex}


\setcounter{ex}{0}
\noindent\textbf{\large PHẦN IV. TỰ LUẬN (3,0 điểm)}

\textit{Thí sinh trình bày chi tiết lời giải vào giấy làm bài.}

\vspace{1mm}
\begin{ex} (1,0 điểm)
Giải hệ phương trình và bất phương trình sau:
\begin{enumerate}[label=\textbf{\alph*)}, leftmargin=0.8cm, itemsep=2pt]
\item (0,5 điểm) $\begin{cases} 2x + y = 1 \\ x - y = 5 \end{cases}$
\item (0,5 điểm) $\frac{x - 3}{5} < \frac{2x + 17}{3}$
\end{enumerate}
\end{ex}
\vspace{0.5mm}

\begin{ex} (0,75 điểm)
Một trường THCS mua $350$ quyển vở (gồm loại $80$ trang và loại $100$ trang) để làm phần thưởng cho học sinh. Giá bán mỗi quyển vở $80$ trang và $100$ trang lần lượt là $6\,000\text{ đồng}$ và $7\,500\text{ đồng}$. Tổng số tiền nhà trường dùng để mua vở là $2\,400\,000\text{ đồng}$. Hỏi nhà trường đã mua mỗi loại bao nhiêu quyển vở?
\end{ex}
\vspace{0.5mm}

\begin{ex} (0,75 điểm)
Lúc 6 giờ sáng, bạn An đi từ nhà ($A$) đến trường ($B$) phải leo lên và xuống một con dốc đỉnh $C$. Cho biết đoạn $AB = 762\text{ m}$, $\widehat{A} = 4^\circ$, $\widehat{B} = 6^\circ$.
\begin{enumerate}[label=\alph*)]
  \item Tính chiều cao con dốc (làm tròn đến hàng đơn vị mét).
  \item Hỏi bạn An đến trường lúc mấy giờ (làm tròn đến phút)? Biết tốc độ lên dốc là $4\text{ km/h}$ và tốc độ xuống dốc là $19\text{ km/h}$.
\end{enumerate}
\begin{center}
\begin{tikzpicture}[scale=1.0, line join=round, line cap=round, >=stealth]
  \coordinate (A) at (0,0); % Nha
  \coordinate (B) at (7.62,0); % Truong
  \coordinate (H) at (4.58,0); % Chan duong cao
  \coordinate (C) at (4.58,1.6); % Dinh doc
  \draw[thick] (-0.5,0) -- (8.2,0) node[below] {Mặt phẳng ngang $(AB)$};
  \draw[line width=1.5pt, blue!75!black] (A) -- (C) -- (B);
  \draw[thick, dashed, red] (C) -- (H) node[midway, right] {$h$};
  \draw (4.58,0.2) -- (4.38,0.2) -- (4.38,0);
  \draw[->, thick, purple] (1.2,0) arc (0:19:1.2);
  \node at (1.5,0.25) [purple] {$4^\circ$};
  \draw[->, thick, purple] (6.42,0) arc (180:152:1.2);
  \node at (6.1,0.25) [purple] {$6^\circ$};
  \filldraw[black] (A) circle (1.5pt) node[below left] {$A$ (Nhà)};
  \filldraw[black] (B) circle (1.5pt) node[below right] {$B$ (Trường)};
  \filldraw[black] (C) circle (1.5pt) node[above] {$C$ (Đỉnh dốc)};
  \filldraw[black] (H) circle (1.5pt) node[below] {$H$};
  \node at (3.8,-0.4) {$AB = 762\text{ m}$};
\end{tikzpicture}
\end{center}
\end{ex}
\vspace{0.5mm}

\begin{ex} (0,5 điểm)
Giải phương trình: $(x + 8)(x + 6)(x + 7)^2 = 72$.
\end{ex}
\vspace{0.5mm}


\enlargethispage{1.5cm}
\begin{center}
\rule{4cm}{0.4pt}\\
\textit{--- HẾT ---}
\end{center}
\end{document}
